\section{Pendahuluan}

\subsection{Latar Belakang}
Perancangan perangkat lunak pada sistem tertanam (\textit{embedded systems}) menuntut kemampuan mikrokontroler untuk merespons kejadian eksternal secara cepat, menjalankan tugas-tugas periodik dengan ketepatan waktu tinggi, serta menjaga kestabilan operasi secara berkelanjutan. Dalam penerapannya pada papan mikrokontroler seperti ESP32, program sederhana yang hanya mengandalkan perulangan utama tunggal (\texttt{loop()}) sering kali tidak memadai ketika sistem harus menangani berbagai pekerjaan secara bersamaan, seperti membaca sinyal sensor berfrekuensi tinggi, memperbarui tampilan antarmuka lokal pada layar OLED, dan menjalankan perhitungan logika kendali secara waktu nyata (\textit{real-time}).

Permasalahan utama dalam pemrograman sekuensial konvensional muncul dari penggunaan metode pembacaan status secara berulang (\textit{polling}) serta fungsi penundaan waktu yang bersifat memblokir (\textit{blocking}) seperti \texttt{delay()}. Pendekatan \textit{polling} memboroskan siklus kerja prosesor dan berisiko melewatkan perubahan sinyal eksternal yang berlangsung singkat, sedangkan penggunaan \texttt{delay()} atau beban komputasi berat menghentikan seluruh alur eksekusi program sehingga tugas periodik mengalami pergeseran waktu dan tampilan antarmuka menjadi tersendat. Selain itu, sistem tertanam yang beroperasi terus-menerus rentan mengalami kondisi macet (\textit{hang}) akibat perulangan tak terbatas atau kegagalan respons perangkat, sehingga diperlukan mekanisme pemulihan baik secara manual maupun otomatis agar sistem tidak berhenti beroperasi tanpa batas waktu.

Berdasarkan permasalahan tersebut, penguasaan mekanisme pengendalian eksekusi tingkat lanjut menjadi kompetensi yang sangat mendesak dalam pembelajaran sistem tertanam. Melalui praktikum Modul 4 ini, mahasiswa mempelajari dan mengimplementasikan lima konsep inti pada mikrokontroler ESP32: pemanfaatan \textit{external interrupt} sebagai mekanisme pemulihan darurat manual (\textit{manual recovery}) saat program utama mengalami \textit{hang}, pembacaan sinyal \textit{quadrature encoder} dua kanal untuk mendeteksi jumlah pulsa dan arah putaran secara presisi, penggunaan \textit{hardware timer interrupt} untuk menjalankan tugas periodik tanpa memblokir program utama, penerapan \textit{watchdog timer} (WDT) untuk mendeteksi dan memulihkan sistem macet secara otomatis, serta implementasi \textit{multitasking} dan komunikasi antar-\textit{task} melalui \textit{queue} menggunakan FreeRTOS pada arsitektur prosesor \textit{dual-core}.

Konsep-konsep yang dipelajari pada modul ini memiliki keterkaitan erat dengan berbagai aplikasi teknologi dunia nyata saat ini. Mekanisme \textit{external interrupt} dan pembacaan \textit{quadrature encoder} merupakan komponen fundamental dalam sistem kendali posisi dan kecepatan motor pada robotika, lengan manipulator industri, serta kendaraan otonom. Di sisi lain, \textit{timer interrupt}, \textit{watchdog timer}, dan penjadwalan tugas paralel menggunakan FreeRTOS menjadi standar arsitektur pada perangkat \textit{Internet of Things} (IoT), sistem akuisisi data medis maupun industri yang membutuhkan determinisme waktu, serta perangkat otonom di lapangan yang wajib pulih secara mandiri ketika terjadi gangguan perangkat lunak.

\subsection{Dasar Teori}

\subsubsection{Ringkasan Mekanisme Pengendalian Eksekusi pada ESP32}
Untuk mengatasi keterbatasan eksekusi sekuensial di dalam fungsi \texttt{loop()}, mikrokontroler ESP32 menyediakan kombinasi mekanisme perangkat keras dan sistem operasi waktu nyata yang saling melengkapi. Karakteristik utama keempat mekanisme yang dibahas dalam modul praktikum ini dirangkum pada Tabel~\ref{tab:ringkasan_modul4}.

\begin{table}[h!]
    \centering
    \caption{Ringkasan Mekanisme Pengendalian Eksekusi dan Keandalan pada Modul 4}
    \label{tab:ringkasan_modul4}
    \vspace{4pt}
    \begin{tabularx}{\textwidth}{>{\bfseries}l X X}
        \toprule
        Mekanisme & Sumber Pemicu & Fungsi Utama pada Sistem \\
        \midrule
        External Interrupt & Perubahan logika sinyal pada pin GPIO eksternal (\texttt{RISING}, \texttt{FALLING}, \texttt{CHANGE}) & Merespons kejadian eksternal secara instan, menghitung pulsa \textit{encoder}, dan pemulihan darurat manual \\
        \addlinespace[4pt]
        Timer Interrupt & Pencacah \textit{hardware timer} internal berdasarkan sinyal \textit{clock} sistem & Menjalankan tugas periodik secara presisi dan \textit{non-blocking} tanpa fungsi \texttt{delay()} \\
        \addlinespace[4pt]
        Watchdog Timer (WDT) & Pencacah waktu pengawas internal yang tidak di-reset dalam batas \textit{timeout} & Mendeteksi kondisi program macet (\textit{hang}) dan memicu \textit{restart} sistem secara otomatis \\
        \addlinespace[4pt]
        FreeRTOS Multitasking & Penjadwal (\textit{scheduler}) berbasis prioritas dan pembagian inti CPU (\textit{dual-core}) & Menjalankan beberapa \textit{task} secara paralel serta bertukar data secara aman melalui \textit{queue} \\
        \bottomrule
    \end{tabularx}
\end{table}

\subsubsection{\textit{External Interrupt} dan \textit{Decoding Quadrature Encoder}}
Pendekatan \textit{polling} memeriksa status logika suatu pin masukan secara terus-menerus di dalam \texttt{loop()}, sehingga kurang efisien untuk kejadian yang jarang muncul atau memerlukan waktu tanggap yang cepat. Sebaliknya, \textbf{\textit{external interrupt}} memungkinkan unit pemroses pusat (CPU) merespons perubahan sinyal pada pin GPIO secara langsung pada tingkat perangkat keras. Ketika kondisi pemicu terpenuhi—baik transisi naik (\texttt{RISING} dari \texttt{LOW} ke \texttt{HIGH}), transisi turun (\texttt{FALLING} dari \texttt{HIGH} ke \texttt{LOW}), maupun setiap perubahan logika (\texttt{CHANGE})—eksekusi program utama dijeda sementara untuk menjalankan fungsi khusus yang disebut \textbf{ISR (\textit{Interrupt Service Routine})}, kemudian CPU kembali melanjutkan program utama dari titik terakhirnya.

Dalam implementasi \textit{external interrupt} pada mikrokontroler ESP32, terdapat beberapa prinsip teknis yang wajib diperhatikan:
\begin{itemize}
    \item \textbf{Keringkasan Fungsi ISR:} Kode di dalam ISR harus dirancang sesingkat dan sesederhana mungkin serta tidak boleh memuat operasi lambat atau memblokir seperti \texttt{delay()} maupun \texttt{Serial.print()} dalam jumlah banyak. Pengecualian berlaku pada skenario tombol darurat (\textit{manual recovery}), di mana fungsi \texttt{esp\_restart()} dipanggil langsung di dalam ISR agar pemulihan tetap berjalan meskipun \texttt{loop()} terjebak dalam kondisi \textit{hang} (\texttt{while (true)}).
    \item \textbf{Kualifikasi Variabel \texttt{volatile}:} Variabel global yang diakses bersama oleh ISR dan program utama harus dideklarasikan menggunakan kata kunci \texttt{volatile} agar kompiler tidak melakukan optimisasi penyimpanan pada register yang menyebabkan pembaruan nilai dari ISR tidak terbaca di \texttt{loop()}.
    \item \textbf{Atribut Memori \texttt{IRAM\_ATTR}:} Fungsi ISR pada ESP32 wajib ditandai dengan atribut \texttt{IRAM\_ATTR} agar instruksinya ditempatkan dan dieksekusi dari memori RAM internal (\textit{Instruction RAM}), bukan dari memori Flash eksternal yang dapat sedang diakses oleh proses lain saat interupsi terjadi.
    \item \textbf{Proteksi Variabel Bersama (\texttt{portMUX\_TYPE}):} Karena ESP32 memiliki dua inti prosesor (\textit{dual-core}), pasangan fungsi \texttt{noInterrupts()} dan \texttt{interrupts()} yang hanya menonaktifkan interupsi pada satu inti tidak cukup aman untuk mencegah konflik akses data. Proteksi yang tepat dilakukan melalui mekanisme \textit{critical section} berbasis \texttt{portMUX\_TYPE} (diinisialisasi dengan \texttt{portMUX\_INITIALIZER\_UNLOCKED}) menggunakan pasangan makro \texttt{portENTER\_CRITICAL\_ISR()} dan \texttt{portEXIT\_CRITICAL\_ISR()} di dalam ISR serta \texttt{portENTER\_CRITICAL()} dan \texttt{portEXIT\_CRITICAL()} di dalam \texttt{loop()}.
\end{itemize}

Salah satu penerapan utama \textit{external interrupt} adalah pembacaan \textbf{\textit{quadrature encoder}} dua kanal (\textit{Channel A} dan \textit{Channel B}) pada motor DC. Dengan memasang interupsi pemicu \texttt{RISING} pada \textit{Channel A} (\texttt{GPIO 32}) dan membaca status logika \textit{Channel B} (\texttt{GPIO 33}) di dalam ISR, perbedaan fase di antara kedua sinyal dapat diterjemahkan untuk menentukan jumlah pulsa sekaligus arah putaran poros motor: apabila \textit{Channel B} berlogika tinggi (\texttt{HIGH}) saat sisi naik \textit{Channel A} terjadi, pencacah pulsa bertambah untuk arah maju (\textit{Clockwise}/CW), sedangkan apabila berlogika rendah (\texttt{LOW}), pencacah pulsa berkurang untuk arah mundur (\textit{Counter-Clockwise}/CCW).

\subsubsection{\textit{Timer Interrupt}}
Selain interupsi yang dipicu oleh perubahan sinyal pin eksternal, ESP32 memiliki empat modul \textbf{\textit{hardware timer}} internal (indeks 0 hingga 3) yang mampu membangkitkan interupsi secara \textbf{periodik} berdasarkan hitungan pulsa waktu internal. Pada konfigurasi timer ESP32 melalui fungsi \texttt{timerBegin(0, 80, true)}, sinyal \textit{clock} dasar sebesar $80\text{ MHz}$ dibagi dengan nilai \textit{prescaler} $80$ sehingga menghasilkan frekuensi pencacahan $1\text{ MHz}$ (setara dengan durasi $1\,\mu\text{s}$ per \textit{tick}) dalam mode hitung naik (\textit{count up}). Nilai batas alarm kemudian ditetapkan melalui \texttt{timerAlarmWrite()} dengan fitur pengisian ulang otomatis (\textit{autoreload}), misalnya sebesar $500.000\text{ tick}$ untuk membangkitkan interupsi tepat setiap $500\text{ ms}$.

Keunggulan utama \textit{timer interrupt} dibandingkan penggunaan fungsi \texttt{delay()} terletak pada sifatnya yang \textbf{\textit{non-blocking}}. Penggunaan \texttt{delay()} menghentikan seluruh eksekusi program selama periode tunggu sehingga interval siklus aktual ikut bergeser apabila terdapat beban komputasi tambahan di dalam \texttt{loop()}. Sebaliknya, \textit{timer interrupt} berjalan secara independen di tingkat perangkat keras dan hanya menyetel variabel penanda (\textit{flag}) di dalam ISR, sementara eksekusi perubahan status keluaran (seperti \textit{toggle} LED pada \texttt{GPIO 2}) dilakukan di dalam \texttt{loop()} tanpa mengganggu alur kerja program utama.

\subsubsection{\textit{Watchdog Timer} (WDT)}
\textbf{\textit{Watchdog Timer} (WDT)} adalah sirkuit pewaktu pengawas perangkat keras yang dirancang untuk menjaga keandalan (\textit{reliability}) sistem tertanam dengan cara memicu \textit{reset} sistem secara otomatis apabila perangkat lunak gagal mereset ulang pencacah WDT—atau disebut istilah "memberi makan" (\textit{feeding the watchdog})—dalam batas waktu tunggu (\textit{timeout}) yang telah ditentukan.

Pada lingkungan Arduino-ESP32, pengawasan ini diimplementasikan menggunakan \textit{Task Watchdog Timer} melalui pustaka \texttt{esp\_task\_wdt.h}. Fungsi \texttt{esp\_task\_wdt\_init(WDT\_TIMEOUT\_S, true)} menginisialisasi periode batas waktu dalam satuan detik beserta pengaktifan mode \textit{panic/restart} otomatis, fungsi \texttt{esp\_task\_wdt\_add(NULL)} mendaftarkan \textit{task} yang sedang berjalan ke dalam daftar pengawasan, dan fungsi \texttt{esp\_task\_wdt\_reset()} menyegarkan kembali hitungan WDT setiap kali tugas berjalan normal. Apabila sistem terjebak dalam kondisi macet (\textit{hang}) dan gagal mengeksekusi \texttt{esp\_task\_wdt\_reset()} hingga waktu tunggu habis, WDT segera memulai ulang ESP32 secara otomatis. Mekanisme pemulihan otomatis berbasis \textit{timeout} ini saling melengkapi dengan pemulihan manual menggunakan \textit{external interrupt}, yang membutuhkan tindakan pengguna tetapi dapat merespons seketika tanpa menunggu batas waktu habis.

\subsubsection{\textit{Multitasking} dengan FreeRTOS}
Kerangka kerja Arduino-ESP32 dibangun di atas sistem operasi waktu nyata \textbf{FreeRTOS} (\textit{Real-Time Operating System}), yang memungkinkan aplikasi dibagi menjadi beberapa unit eksekusi independen yang disebut \textbf{\textit{task}}. Setiap \textit{task} memiliki fungsi perulangan, ukuran memori tumpukan (\textit{stack}), dan tingkat prioritas tersendiri. Pada mikrokontroler ESP32 yang memiliki dua inti pemroses (\textit{Core 0} dan \textit{Core 1}), fungsi \texttt{xTaskCreatePinnedToCore()} memungkinkan setiap \textit{task} dijalankan secara paralel sejati pada inti prosesor yang berbeda, sementara fungsi \texttt{loop()} utama dapat dibiarkan kosong karena penjadwalan diatur sepenuhnya oleh \textit{scheduler} FreeRTOS.

Agar perpindahan antar-\textit{task} berlangsung efisien, penundaan waktu di dalam \textit{task} menggunakan fungsi \texttt{vTaskDelay(pdMS\_TO\_TICKS(...))} yang berbasis satuan \textit{tick} FreeRTOS dan bersifat \textit{non-blocking} terhadap \textit{task} lain. Selanjutnya, untuk pertukaran data antar-\textit{task} yang aman dari konflik akses bersamaan (\textit{thread-safe}), FreeRTOS menyediakan mekanisme antrean data yang disebut \textbf{\textit{queue}} (\texttt{QueueHandle\_t}). Melalui pembuatan antrean berkapasitas satu data menggunakan \texttt{xQueueCreate(1, sizeof(double))}, sebuah \textit{task} komputasi berat pada \textit{Core 0} dapat memperbarui hasil perhitungan terkini secara langsung memakai \texttt{xQueueOverwrite()}, sementara \textit{task} pembaruan tampilan layar OLED pada \textit{Core 1} mengambil data tersebut tanpa menunggu menggunakan \texttt{xQueueReceive()} dengan parameter waktu tunggu nol, sehingga animasi tampilan tetap berjalan mulus tanpa terhambat proses komputasi.

%===========================================================%
\section{Tugas Pendahuluan}
\begin{enumerate}[label=\textbf{\arabic*.}, leftmargin=*, itemsep=10pt]

    \item \textbf{Perbedaan mendasar antara \textit{polling} dan \textit{interrupt} dalam menangani kejadian pada mikrokontroler}

    Perbedaan mendasar antara metode \textit{polling} dan \textit{interrupt} terletak pada cara prosesor (CPU) mendeteksi serta merespons perubahan kondisi eksternal, seperti penekanan tombol pada pin masukan mikrokontroler. Perbandingan karakteristik utama kedua metode tersebut dirangkum pada Tabel~\ref{tab:polling_vs_interrupt}.

    \begin{table}[h!]
        \centering
        \caption{Perbandingan Metode \textit{Polling} dan \textit{Interrupt} pada Mikrokontroler}
        \label{tab:polling_vs_interrupt}
        \vspace{4pt}
        \begin{tabularx}{\textwidth}{>{\bfseries}l X X}
            \toprule
            Aspek Pembanding & Polling & Interrupt \\
            \midrule
            Mekanisme Deteksi & CPU memeriksa status pin secara berulang dan aktif di dalam program utama & Perangkat keras memantau perubahan sinyal dan mengirim sinyal interupsi ke CPU \\
            \addlinespace[4pt]
            Pemanfaatan CPU & Kurang efisien karena siklus CPU terpakai terus-menerus meski tidak ada kejadian & Sangat efisien karena CPU bebas menjalankan tugas lain hingga kejadian muncul \\
            \addlinespace[4pt]
            Waktu Respons & Bergantung pada lama eksekusi satu siklus perulangan utama (\textit{loop}) & Sangat cepat dan deterministik segera setelah transisi sinyal terdeteksi \\
            \addlinespace[4pt]
            Risiko Kehilangan Sinyal & Tinggi apabila penekanan tombol berlangsung singkat saat CPU sibuk & Sangat rendah karena transisi sinyal langsung memicu perangkat keras \\
            \bottomrule
        \end{tabularx}
    \end{table}

    \begin{itemize}
        \item \textbf{Metode \textit{Polling}:} Program utama secara berkala membaca status logika pin tombol (misalnya melalui pemanggilan fungsi \texttt{digitalRead()}) di dalam perulangan terus-menerus. Apabila tombol ditekan dan dilepas secara cepat ketika CPU sedang mengeksekusi instruksi lain atau fungsi tunda, perubahan sinyal tersebut dapat terlewatkan. Selain itu, waktu prosesor terbuang untuk memeriksa tombol meskipun tidak terjadi perubahan status.
        \item \textbf{Metode \textit{Interrupt}:} Rangkaian perangkat keras pengendali GPIO pada mikrokontroler dikonfigurasi untuk mendeteksi transisi tepi sinyal tertentu (\textit{rising edge}, \textit{falling edge}, atau \textit{change}). Saat tombol ditekan dan memicu transisi logika, perangkat keras segera membangkitkan permintaan interupsi (\textit{Interrupt Request}). CPU menghentikan sementara eksekusi program utama, menyimpan konteks register saat ini, mengeksekusi fungsi penanganan khusus yang disebut \textit{Interrupt Service Routine} (ISR), lalu kembali melanjutkan instruksi yang tertunda.
    \end{itemize}

    \item \textbf{Alasan kode di dalam ISR harus singkat dan sederhana serta risiko jika berjalan terlalu lama}

    Kode di dalam \textit{Interrupt Service Routine} (ISR) wajib dirancang sesingkat dan sesederhana mungkin karena selama ISR dieksekusi, alur eksekusi program utama serta interupsi lain dengan prioritas yang sama atau lebih rendah ditangguhkan sementara. Praktik terbaik dalam perancangan ISR adalah membatasi tugasnya hanya pada operasi cepat yang bersifat atomik, seperti mengubah nilai variabel penanda (\textit{flag}) bertipe \texttt{volatile}, menambah nilai pencacah (\textit{counter}), mencatat cap waktu (\textit{timestamp}), atau mengirim data ringkas ke antrean menggunakan fungsi khusus ISR (seperti \texttt{xQueueSendFromISR()}), sementara pemrosesan data yang berat dialihkan ke perulangan utama atau \textit{task} terpisah.

    Apabila kode di dalam ISR berjalan terlalu lama, sistem berisiko mengalami beberapa gangguan serius berikut:
    \begin{itemize}
        \item \textbf{Kehilangan Interupsi Berikutnya (\textit{Missed Interrupts}):} Interupsi dari sumber periferal lain yang datang saat CPU masih tertahan di dalam ISR dapat tertunda atau hilang apabila terjadi lebih dari satu kali sebelum interupsi sebelumnya selesai dilayani.
        \item \textbf{Menurunnya Responsivitas Sistem (\textit{High Latency}):} Tugas-tugas pada program utama maupun periferal komunikasi waktu nyata mengalami keterlambatan eksekusi sehingga perilaku sistem menjadi tidak stabil atau mengalami pergeseran waktu (\textit{jitter}).
        \item \textbf{Pemicuan \textit{Reset} oleh \textit{Interrupt Watchdog}:} Pada mikrokontroler seperti ESP32, terdapat pengawas khusus (\textit{Interrupt Watchdog Timer}) yang memantau lamanya waktu penonaktifan interupsi; jika sebuah ISR menahan CPU melebihi batas waktu tersebut, sistem dianggap macet dan memicu \textit{reset} perangkat keras secara paksa.
        \item \textbf{Kegagalan Sistem atau \textit{Deadlock}:} Pemanggilan fungsi yang bersifat memblokir (\textit{blocking}), mencetak pesan panjang ke port serial, atau mengalokasikan memori dinamis di dalam ISR dapat menyebabkan sistem terkunci permanen karena mekanisme penjadwalan dan interupsi pendukungnya sedang tidak aktif.
    \end{itemize}

    \item \textbf{Perbedaan antara \textit{timer interrupt} dan penggunaan fungsi \texttt{delay()} untuk tugas periodik}

    Meskipun keduanya dapat digunakan untuk mengatur jeda waktu eksekusi suatu tugas, \textit{timer interrupt} dan fungsi \texttt{delay()} memiliki mekanisme kerja serta dampak kinerja yang sangat berbeda:
    \begin{itemize}
        \item \textbf{Sifat Eksekusi (\textit{Non-blocking} vs \textit{Blocking}):} Fungsi \texttt{delay()} bekerja secara memblokir (\textit{blocking}) dengan menahan alur eksekusi CPU pada baris kode tersebut hingga durasi waktu tunda berakhir, sehingga CPU tidak dapat menjalankan instruksi lain pada alur yang sama. Sebaliknya, \textit{timer interrupt} memanfaatkan modul pencacah perangkat keras (\textit{hardware timer}) yang berjalan secara mandiri di latar belakang berdasarkan sinyal \textit{clock} sistem, sehingga CPU tetap bebas mengeksekusi tugas-tugas lain selama jeda waktu berlangsung.
        \item \textbf{Presisi dan Akurasi Periode:} Pada penggunaan fungsi \texttt{delay()}, periode pengulangan tugas sebenarnya adalah penjumlahan antara waktu tunda dan lama eksekusi kode tugas itu sendiri ($T_{\text{total}} = T_{\text{delay}} + T_{\text{eksekusi}}$), sehingga interval waktu menjadi bergeser (\textit{drift}) apabila beban komputasi berubah. Sebaliknya, \textit{timer interrupt} membangkitkan interupsi tepat setiap kali pencacah perangkat keras mencapai nilai batasnya secara konsisten, menghasilkan interval periodik yang jauh lebih presisi dan deterministik.
        \item \textbf{Skalabilitas Sistem:} Penggunaan banyak fungsi \texttt{delay()} menyulitkan pengelolaan beberapa tugas periodik dengan interval berbeda secara bersamaan, sedangkan \textit{timer interrupt} memungkinkan sistem menjadwalkan berbagai aktivitas periodik tanpa saling menghambat.
    \end{itemize}

    \item \textbf{Cara kerja \textit{watchdog timer} dalam mendeteksi dan memulihkan sistem dari kondisi \textit{hang}}

    \textit{Watchdog Timer} (WDT) adalah sirkuit pencacah mundur perangkat keras (\textit{hardware countdown timer}) yang beroperasi secara independen dari alur eksekusi utama CPU sebagai mekanisme pengaman otomatis terhadap kegagalan perangkat lunak. Tahapan kerja WDT dalam mendeteksi dan memulihkan kondisi sistem yang macet (\textit{hang}) berlangsung sebagai berikut:
    \begin{enumerate}[label=\alph*., leftmargin=*]
        \item \textbf{Inisialisasi dan Pencacahan Mundur:} Ketika sistem mulai berjalan, WDT diaktifkan dengan nilai batas waktu tunggu tertentu (\textit{timeout period}) dan mulai menghitung mundur secara terus-menerus menggunakan sumber sinyal pewaktuan tersendiri.
        \item \textbf{Penyegaran Berkala (\textit{Feeding / Kicking the Watchdog}):} Selama program berjalan normal tanpa hambatan, perangkat lunak secara berkala menyetel ulang nilai pencacah WDT kembali ke nilai awalnya sebelum hitungan mencapai nol (misalnya melalui pemanggilan fungsi \texttt{esp\_task\_wdt\_reset()} pada ESP32).
        \item \textbf{Deteksi Kondisi Macet (\textit{Timeout}):} Apabila perangkat lunak terjebak dalam perulangan tak terbatas (\textit{infinite loop}), mengalami kebuntuan akses sumber daya (\textit{deadlock}), atau terhenti akibat gangguan elektromagnetik, CPU gagal menjalankan instruksi penyegaran WDT. Akibatnya, pencacah mundur WDT terus berkurang hingga mencapai nol.
        \item \textbf{Pemulihan Otomatis (\textit{System Reset}):} Sesaat setelah pencacah WDT mencapai nol, perangkat keras WDT segera membangkitkan sinyal \textit{reset} sistem untuk memulai ulang mikrokontroler dari kondisi awal yang stabil tanpa memerlukan campur tangan operator secara manual.
    \end{enumerate}

    \item \textbf{Konsep \textit{task} dan \textit{queue} pada FreeRTOS serta kegunaannya pada mikrokontroler \textit{dual-core} seperti ESP32}

    \begin{itemize}
        \item \textbf{Konsep \textit{Task}:} Pada sistem operasi waktu nyata FreeRTOS, \textit{task} merupakan unit program independen yang menjalankan fungsi tertentu di dalam perulangannya sendiri. Setiap \textit{task} memiliki alokasi memori tumpukan (\textit{stack}), konteks register, serta tingkat prioritas tersendiri. Penjadwal (\textit{scheduler}) FreeRTOS mengatur pergantian eksekusi antar-\textit{task} berdasarkan prioritas dan statusnya (\textit{Running}, \textit{Ready}, \textit{Blocked}, atau \textit{Suspended}) sehingga berbagai proses dapat berjalan secara terstruktur.
        \item \textbf{Konsep \textit{Queue}:} \textit{Queue} adalah struktur penyangga antrean berbasis prinsip FIFO (\textit{First-In, First-Out}) yang aman dari konflik akses bersamaan (\textit{thread-safe}) untuk pertukaran data antar-\textit{task} maupun dari ISR menuju \textit{task}. Data yang dikirim ke dalam \textit{queue} (misalnya melalui \texttt{xQueueSend()}) disalin ke ruang memori antrean, dan apabila antrean kosong, \textit{task} penerima (\texttt{xQueueReceive()}) dapat masuk ke kondisi \textit{Blocked} tanpa membebani siklus CPU hingga data baru masuk.
        \item \textbf{Kegunaan pada Mikrokontroler \textit{Dual-Core} (ESP32):} Mikrokontroler ESP32 memiliki dua inti pemroses terpisah (\textit{Core 0} dan \textit{Core 1}). Melalui pembagian program menjadi beberapa \textit{task} (misalnya menggunakan fungsi \texttt{xTaskCreatePinnedToCore()}), beban kerja sistem dapat didistribusikan secara paralel sejati di antara kedua inti CPU—seperti menugaskan \textit{Core 0} untuk menangani komunikasi jaringan atau Wi-Fi sementara \textit{Core 1} berfokus pada akuisisi sensor dan kendali aktuator waktu nyata. Dalam arsitektur dua inti tersebut, \textit{queue} berperan penting sebagai jembatan pertukaran data yang aman sehingga mencegah kerusakan data (\textit{race condition}) ketika dua \textit{task} pada inti berbeda mengakses informasi yang sama.
    \end{itemize}

    \item \textbf{Analogi sederhana dalam kehidupan sehari-hari terkait mekanisme WDT, ISR, dan FreeRTOS}

    Ketiga mekanisme tersebut dapat dianalogikan melalui aktivitas kerja operasional pada sebuah dapur restoran yang sibuk:
    \begin{itemize}
        \item \textbf{Analogi ISR (\textit{Interrupt Service Routine}):} Seorang koki sedang fokus memasak hidangan utama di dapur tanpa harus berjalan ke pintu depan setiap beberapa detik untuk memeriksa apakah kurir bahan datang (\textit{polling}). Ketika bel pintu depan berbunyi (\textit{interrupt}), koki segera mengecilkan api kompor, membuka pintu sebentar untuk menerima dan meletakkan paket bahan di meja penerimaan (\textit{ISR} yang singkat), lalu langsung kembali mengaduk masakannya sebelum membuka rincian paket tersebut nanti.
        \item \textbf{Analogi WDT (\textit{Watchdog Timer}):} Di dapur tersebut terpasang alat pengingat otomatis yang berbunyi setiap 5 menit; setiap kali koki selesai mengecek kondisi kompor dan menekan tombol konfirmasi (\textit{feeding the watchdog}), waktu pengingat kembali diatur ulang ke 5 menit. Apabila koki tiba-tiba pingsan atau terjebak masalah hingga tidak menekan tombol konfirmasi sampai waktu habis, sistem pengaman dapur secara otomatis memutus aliran gas dan menyalakan sistem pemulihan darurat (\textit{system reset}) guna mencegah kebakaran.
        \item \textbf{Analogi FreeRTOS (\textit{Task} dan \textit{Queue}):} Dapur restoran tersebut dikelola oleh seorang kepala dapur (\textit{scheduler}) yang membagi pekerjaan ke dua meja kerja utama (\textit{dual-core}) dengan peran terpisah (\textit{task}), yaitu bagian persiapan bahan, bagian pemasakan sesuai tingkat urgensi pesanan (\textit{prioritas}), dan bagian penyajian. Agar koki persiapan dan koki pemasakan tidak saling berebut talenan atau berteriak bersamaan, keduanya menggunakan **rak nampan pesanan berurutan (\textit{queue})**: bahan yang telah dipotong diletakkan berurutan ke nampan, lalu diambil oleh koki pemasakan sesuai urutan kedatangannya.
    \end{itemize}

\end{enumerate}